\section{\uppercase{Discussion}}
\label{sec:discussion}

\noindent
The three experiments indicate that the classifiers use the pen dynamics encoded in the minimaps. When the static trace was available and favoured by ties, training-only selection chose a dynamic encoding in 24 of 25 training sets in Experiment~2, and in 23 of 25 with LPQ and 24 of 25 with ResNet18 in Experiment~3. Displacing the signals within strokes lowered recording-level macro F1, most under joint permutation, so the position of each value along the trajectory matters and not only the colours present. The benefit is clearest on the spiral, the drawing task, where $SAZ$ added 0.0802 in Experiment~1 and permutation removed 0.0609 in Experiment~2.

The gains depend on the task (Figure~\ref{fig:heatmap}). One encoding chosen for all eight tasks averages gains on the spiral and T6 with losses on T2, T3 and T5, which keeps its eight-task mean close to the static trace, whereas on the homogeneous letter patterns of T2--T4 the selected encoding gives the best participant-level decisions. Choosing encodings per task, rather than one for the whole protocol, is a direct way to exploit this dependence.

The pipelines are strong on different tasks (Table~\ref{tab:per-task}): minimaps with LPQ lead the spiral, kinematic attributes lead T2, T3 and the sentence, and minimaps with ResNet18 lead T4--T7 and the eight-task mean. The representation therefore works with both a handcrafted texture descriptor and pretrained convolutional features. The complementarity between the image and kinematic pipelines, each strongest on a different endpoint, motivates combining them. Shortening the protocol from eight tasks to three improved both the minimap LPQ and the kinematic pipelines, which reduces the number of exercises, although assessment time and clinical burden were not evaluated.

\section{\uppercase{Conclusion}}
\label{sec:conclusion}
\noindent
This initial study introduced pseudodynamic minimaps, a new image representation that places pen speed, pressure, altitude and azimuth as colour along the handwriting trajectory. Evaluated on PaHaW with participant-separated partitions, the representation was preferred over static traces by training-only selection, displacing the dynamics along each stroke lowered its performance, and on the three repeated-pattern tasks it gave the best participant-level decisions among the four compared pipelines. These results indicate that pseudodynamic minimaps are an effective way to bring pen dynamics into image-based handwriting analysis for Parkinson's disease.

The evidence comes from 75 participants of a single dataset that also informed the design of the experiments, so the size of the gains remains uncertain and independent cohorts are needed to confirm them. A source package with the code and configurations of all experiments supports that validation; the original recordings are not redistributed.

Future work will choose encodings per task, combine minimaps with kinematic descriptors, and fine-tune pretrained networks on the minimaps. A direct comparison with published PaHaW results is also left for future work. It requires more than citing reported accuracies: each method has to be re-implemented and evaluated under the same participant-separated partitions, without selecting tasks or features on test participants.
